\documentclass[11pt,letterpaper]{article}
% preamble.tex — Shared preamble for all documents
\usepackage{graphicx}
\usepackage{hyperref}
\usepackage{enumitem}
\usepackage{titlesec}
\usepackage{fancyhdr}
\usepackage{setspace}
\usepackage{parskip}
\usepackage{etoolbox}
\usepackage{xspace}
\usepackage{needspace}
\usepackage{changepage}
\usepackage{geometry}
\usepackage{ragged2e}
\usepackage{environ}

% Hyperref setup
\hypersetup{
  colorlinks=false,
  pdfborder={0 0 0},
  pdftitle={The Letter to the Father},
  pdfauthor={},
}

% Load project files
% colors.tex — Color palette for The Letter to the Father
\usepackage{xcolor}

% Primary palette — muted, period-appropriate, prestige
\definecolor{kafkablack}{HTML}{1A1A1A}
\definecolor{kafkagray}{HTML}{4A4A4A}
\definecolor{kafkalight}{HTML}{8A8A8A}
\definecolor{kafkacream}{HTML}{F5F0E8}
\definecolor{kafkasepia}{HTML}{704214}
\definecolor{kafkablood}{HTML}{6B1C23}
\definecolor{kafkaink}{HTML}{2C2C3A}

% Accent colors — for package materials
\definecolor{prestigegold}{HTML}{B8860B}
\definecolor{sectionrule}{HTML}{3A3A3A}
\definecolor{pullquote}{HTML}{5A4A3A}

% metadata.tex — Project metadata for The Letter to the Father
\newcommand{\filmtitle}{The Letter to the Father}
\newcommand{\filmtitleupper}{THE LETTER TO THE FATHER}
\newcommand{\filmsubtitle}{An Original Screenplay}
\newcommand{\filmauthor}{Based on the life and writings of Franz Kafka}
\newcommand{\filmyear}{2026}
\newcommand{\filmdraft}{Revised Draft -- Public Development Example}
\newcommand{\filmdate}{October 2026 -- Example Edition}
\newcommand{\filmlogline}{In the final years of his life, a frail and tormented Franz Kafka struggles to write, love, and exist within a world of family pressure, bureaucratic humiliation, illness, and spiritual dread --- while the stories he creates begin to mirror and consume the life he can no longer control.}
\newcommand{\filmgenre}{Prestige Drama / Literary Biopic}
\newcommand{\filmrating}{R}
\newcommand{\filmlead}{Adrien Brody}
\newcommand{\filmbudgetrange}{\$12--18M}
\newcommand{\filmruntime}{Approximately 110 minutes}


% Public example status travels with every downloadable document.
\AddToHook{shipout/foreground}{\put(40,-776){\makebox(0,0)[l]{\fontsize{8}{9}\selectfont\textit{Development example. All casting, financing and awards language is illustrative; no attachments confirmed.}}}}

% typography.tex — Typography settings
\usepackage{fontspec}
\usepackage{microtype}

% Prefer repo-local metric-compatible fonts so Tectonic does not depend on
% absolute host-system font paths. Fall back to system fonts only if the vendored
% font payload is unavailable.
\IfFileExists{../../../fonts/Tinos-Regular.ttf}{
  \setmainfont{Tinos-Regular.ttf}[
    Path=../../../fonts/,
    BoldFont=Tinos-Bold.ttf,
    ItalicFont=Tinos-Italic.ttf,
    BoldItalicFont=Tinos-BoldItalic.ttf,
  ]
}{
  \IfFontExistsTF{TeX Gyre Termes}{
    \def\projectmainfont{TeX Gyre Termes}
  }{
    \IfFontExistsTF{Times New Roman}{
      \def\projectmainfont{Times New Roman}
    }{
      \def\projectmainfont{Latin Modern Roman}
    }
  }
  \setmainfont{\projectmainfont}
}

\IfFileExists{../../../fonts/Arimo-Regular.ttf}{
  \setsansfont{Arimo-Regular.ttf}[
    Path=../../../fonts/,
    BoldFont=Arimo-Bold.ttf,
    ItalicFont=Arimo-Italic.ttf,
    BoldItalicFont=Arimo-BoldItalic.ttf,
  ]
}{
  \IfFontExistsTF{TeX Gyre Heros}{
    \def\projectsansfont{TeX Gyre Heros}
  }{
    \IfFontExistsTF{Helvetica}{
      \def\projectsansfont{Helvetica}
    }{
      \def\projectsansfont{Latin Modern Sans}
    }
  }
  \setsansfont{\projectsansfont}
}

\IfFileExists{../../../fonts/Cousine-Regular.ttf}{
  \setmonofont{Cousine-Regular.ttf}[
    Path=../../../fonts/,
    BoldFont=Cousine-Bold.ttf,
    ItalicFont=Cousine-Italic.ttf,
    BoldItalicFont=Cousine-BoldItalic.ttf,
    Scale=0.95,
  ]
  \def\screenplayfont{%
    \fontspec{Cousine-Regular.ttf}[%
      Path=../../../fonts/,
      BoldFont=Cousine-Bold.ttf,
      ItalicFont=Cousine-Italic.ttf,
      BoldItalicFont=Cousine-BoldItalic.ttf,
      Scale=1.0,
    ]%
  }
}{
  \IfFontExistsTF{Courier New}{
    \def\projectmonofont{Courier New}
  }{
    \IfFontExistsTF{Courier}{
      \def\projectmonofont{Courier}
    }{
      \def\projectmonofont{Latin Modern Mono}
    }
  }
  \setmonofont{\projectmonofont}[Scale=0.95]
  \def\screenplayfont{\fontspec{\projectmonofont}[Scale=1.0]}
}

% Microtype settings
\microtypesetup{
  protrusion=true,
  expansion=false,
}

% studio_package_macros.tex — Formatting macros for the studio package

% ============================================================
% SECTION STYLING
% ============================================================
\titleformat{\section}
  {\Large\bfseries\color{kafkablack}\sffamily}
  {}
  {0em}
  {}
  [\vspace{2pt}\textcolor{sectionrule}{\rule{\textwidth}{0.5pt}}\vspace{6pt}]

\titleformat{\subsection}
  {\large\bfseries\color{kafkagray}\sffamily}
  {}
  {0em}
  {}

\titleformat{\subsubsection}
  {\normalsize\bfseries\color{kafkagray}\sffamily}
  {}
  {0em}
  {}

% ============================================================
% PULL QUOTE
% ============================================================
% Usage: \pullquote{Text here}
\NewEnviron{pullquote}{%
  \vspace{8pt}%
  \begin{adjustwidth}{0.5in}{0.5in}%
    \color{pullquote}\itshape\large%
    \BODY%
  \end{adjustwidth}%
  \vspace{8pt}%
}

% ============================================================
% HIGHLIGHT BOX
% ============================================================
\newcommand{\highlightbox}[1]{%
  \par\vspace{6pt}%
  \noindent\colorbox{kafkacream}{%
    \parbox{\dimexpr\textwidth-2\fboxsep}{%
      \vspace{4pt}#1\vspace{4pt}%
    }%
  }%
  \par\vspace{6pt}%
}

% ============================================================
% TARGET ENTRY (for director/producer lists)
% ============================================================
% Usage: \target{Name}{Why they fit}{Tier}
\newcommand{\target}[3]{%
  \needspace{4\baselineskip}%
  \subsubsection*{#1 \hfill \textnormal{\small\textit{#3}}}%
  #2\par\vspace{6pt}%
}

% ============================================================
% SCENE CARD (for showcase section)
% ============================================================
% Usage: \scenecard{Title}{Description}
\newcommand{\scenecard}[2]{%
  \needspace{5\baselineskip}%
  \vspace{4pt}%
  \noindent\textbf{\sffamily #1}\\[3pt]%
  #2\par\vspace{8pt}%
  \textcolor{sectionrule}{\rule{0.3\textwidth}{0.25pt}}%
  \vspace{8pt}%
}

% ============================================================
% TAGLINE
% ============================================================
\newcommand{\tagline}[1]{%
  \vspace{4pt}%
  \begin{center}%
    {\large\itshape\color{kafkagray} ``#1''}%
  \end{center}%
  \vspace{4pt}%
}

% ============================================================
% COMP FILM ENTRY
% ============================================================
% Usage: \compfilm{Title (Year)}{Relevance}
\newcommand{\compfilm}[2]{%
  \noindent\textbf{#1}\\%
  {\small #2}\par\vspace{6pt}%
}

% ============================================================
% RISK ENTRY
% ============================================================
% Usage: \risk{Risk}{Mitigation}
\newcommand{\risk}[2]{%
  \noindent\textbf{Risk:} #1\\%
  \textbf{Mitigation:} #2\par\vspace{8pt}%
}

% ============================================================
% PACKAGE GEOMETRY
% ============================================================
\newcommand{\packagegeometry}{%
  \newgeometry{
    left=1in,
    right=1in,
    top=1in,
    bottom=1in,
  }%
}

\geometry{left=0.75in, right=0.75in, top=0.75in, bottom=0.75in}
\pagestyle{empty}
\setlength{\parindent}{0pt}
\setlength{\parskip}{5pt}

\begin{document}

\begin{center}
{\Large\bfseries\sffamily \filmtitleupper}\\[3pt]
{\normalsize\sffamily\color{kafkagray} Festival Programmer Brief}
\end{center}
\vspace{2pt}
\textcolor{sectionrule}{\rule{\textwidth}{0.5pt}}
\vspace{4pt}

\textbf{Why It Belongs in Competition}\\
A prestige feature with a clean festival proposition: a major actor at the height of his powers, a culturally central subject, and a formal approach that is disciplined rather than decorative. The film does not treat Kafka as heritage material. It turns his life into a pressure system: office, family, illness, love, and authorship all closing around the same undelivered sentence. That gives the project both the seriousness programmers want and the dramatic shape audiences actually respond to.

\vspace{2pt}

\textbf{The Programming Value}\\
The film sits at the useful intersection of accessibility and rigor. It is undeniably actor-forward, which gives it immediate public legibility, but it also has genuine formal identity: six movements instead of standard biopic chronology, recurring institutional geometry, and a title object that keeps re-entering the film as contamination rather than explanation. It can play for serious cinephiles, awards press, and prestige audiences without becoming a compromise object for any of them.

\vspace{2pt}

\textbf{Why Venice Is the Cleanest Fit}\\
Venice gives the film the most natural frame: European subject, severe control, emotional clarity, and obvious awards potential through Brody. The project belongs in the lane where a competition title can be both aesthetically exact and strategically useful in the Oscar corridor. Cannes is equally plausible if the final director attachment pushes the package more overtly toward auteur branding; Telluride and TIFF remain the ideal acceleration points after a major premiere.

\vspace{2pt}

\textbf{What the Festival Audience Gets}\\
Not a literary-history lesson. Not a dutiful biographical pageant. The office material gives the first half procedural traction. The father-letter provides a memorable spine. Felice, Milena, and Dora each redirect the emotional temperature of the film. Berlin gives the audience a brief, startling warmth before the final room. The result is a film with severity, but not monotony.

\vspace{2pt}

\textbf{Key Competition Assets}\\
\textit{Performance}: a role built for Brody's exact strengths, with multiple clip scenes and a full physical deterioration arc. --- \textit{Form}: prestige realism bent just enough to feel Kafkaesque without slipping into literary illustration. --- \textit{Place}: Prague, Berlin, Vienna, and the Czech countryside rendered as repeatable emotional systems rather than tourist period décor. --- \textit{Conversation}: a title and premise that can be explained instantly, then deepen on contact.

\vspace{2pt}

\textbf{Post-Premiere Momentum}\\
The film is built to convert a competition launch into sustained awards and specialty-theatrical life. A major festival premiere gives the press story immediate shape: Brody, Kafka, the father-letter, the formal control, the return of adult prestige drama that does not confuse refinement with inertia. It is exactly the kind of title that benefits from curation first and campaign second.

\vspace{4pt}
\textcolor{sectionrule}{\rule{\textwidth}{0.5pt}}
\vspace{2pt}
\begin{center}
{\small\color{kafkalight}\textit{A competition title with real performance gravity, real formal identity, and a cleaner emotional hook than most literary-biographical fare.}}
\end{center}

\end{document}
